%Môn: Toán
\begin{ex}
    Tính giới hạn $L = \lim_{n \to \infty} \frac{1}{n^2}$.
    \shortans{0}
    \loigiai{Khi $n$ dần tới vô cực, $n^2$ cũng dần tới vô cực, do đó $1/n^2$ dần tới 0.}
\end{ex}
%Môn: Vật Lí
\begin{ex}
    Một vật dao động điều hòa với biên độ $A = 4$ cm và chu kì $T = 2$ s. Tìm tần số góc $\omega$ của dao động.
    \shortans{3,14}
    \loigiai{Tần số góc $\omega = 2\pi/T = 2\pi/2 = \pi \approx 3,14$ rad/s.}
\end{ex}
%Môn: Hóa Học
\begin{ex}
    Trong phân tử methane ($CH_4$), nguyên tử carbon có hóa trị là bao nhiêu?
    \shortans{4}
    \loigiai{Trong các hợp chất hữu cơ, nguyên tử carbon luôn có hóa trị 4.}
\end{ex}
%Môn: Địa Lí
\begin{ex}
    Năm 2021, EU có bao nhiêu quốc gia thành viên?
    \shortans{27}
    \loigiai{Theo sách giáo khoa Địa Lí 11 trang 37, năm 2021 Liên minh châu Âu có 27 thành viên sau khi Anh rời đi.}
\end{ex}
%Môn: Lịch Sử
\begin{ex}
    Năm nào vua Minh Mạng thực hiện cuộc cải cách hành chính chia cả nước thành 30 tỉnh và 1 phủ Thừa Thiên?
    \shortans{1831}
    \loigiai{Cuộc cải cách hành chính của Minh Mạng bắt đầu từ năm 1831 (phía Bắc) và hoàn thành vào năm 1832 (phía Nam).}
\end{ex}
%Môn: Giáo Dục Kinh Tế và Pháp Luật
\begin{ex}
    Theo quy định của pháp luật Việt Nam, độ tuổi tối thiểu để công dân có quyền bầu cử là bao nhiêu?
    \shortans{18}
    \loigiai{Công dân đủ 18 tuổi trở lên có quyền bầu cử đại biểu Quốc hội và Hội đồng nhân dân các cấp.}
\end{ex}
%Môn: Toán
\begin{ex}
    Cho cấp số nhân $(u_n)$ có $u_1 = 2$ và công bội $q = 3$. Tính số hạng $u_2$.
    \shortans{6}
    \loigiai{Áp dụng công thức $u_n = u_{n-1} \cdot q$, ta có $u_2 = 2 \cdot 3 = 6$.}
\end{ex}
%Môn: Vật Lí
\begin{ex}
    Tốc độ âm thanh trong không khí ở điều kiện tiêu chuẩn xấp xỉ bao nhiêu m/s?
    \shortans{340}
    \loigiai{Tốc độ truyền âm trong không khí thường được lấy là 340 m/s.}
\end{ex}
%Môn: Hóa Học
\begin{ex}
    Số oxi hóa của nitrogen trong phân tử $NH_3$ là bao nhiêu?
    \shortans{-3}
    \loigiai{Trong $NH_3$, hydrogen có số oxi hóa +1, do đó nitrogen có số oxi hóa -3.}
\end{ex}
%Môn: Địa Lí
\begin{ex}
    Khu vực Mỹ La-tinh có diện tích khoảng bao nhiêu triệu km$^2$?
    \shortans{20}
    \loigiai{Theo sách giáo khoa Địa lí 11 trang 22, khu vực Mỹ La-tinh có diện tích khoảng 20 triệu km$^2$.}
\end{ex}
%Môn: Tiếng Anh
\begin{ex}
    Which of the following is NOT a cause of global warming?
    \choice
    {Burning fossil fuels}
    {Deforestation}
    {\True Planting trees}
    {Industrial waste}
    \loigiai{Planting trees helps absorb $CO_2$, thus reducing global warming, not causing it.}
\end{ex}
%Môn: Vật Lí
\begin{ex}
    \immini[thm]{
        Cho đồ thị li độ - thời gian của một vật dao động điều hòa như hình bên. Hãy xác định chu kì dao động $T$ (tính theo giây).
    }{
        [Image showing one full cycle in 2 seconds]}
    \shortans{2}
    \loigiai{Dựa vào đồ thị, khoảng thời gian để vật thực hiện một dao động toàn phần là 2 giây.}
\end{ex}
%Môn: Hóa Học
\begin{ex}
    Hợp chất $C_2H_5OH$ thuộc loại chất nào sau đây?
    \choice
    {Hydrocarbon}
    {Dẫn xuất halogen}
    {\True Alcohol}
    {Phenol}
    \loigiai{$C_2H_5OH$ (ethanol) có nhóm chức $-OH$ liên kết với gốc alkyl nên là alcohol.}
\end{ex}
%Môn: Toán
\begin{ex}
    Về quan hệ vuông góc trong không gian, khẳng định nào sau đây là sai?
    \choiceTF
    {\True Hai đường thẳng cùng vuông góc với một đường thẳng thứ ba thì song song với nhau.}
    {Một đường thẳng vuông góc với một trong hai mặt phẳng song song thì cũng vuông góc với mặt phẳng kia.}
    {Hai mặt phẳng phân biệt cùng vuông góc với một đường thẳng thì song song với nhau.}
    {Đường thẳng vuông góc với mặt phẳng thì vuông góc với mọi đường thẳng nằm trong mặt phẳng đó.}
    \loigiai{
        \begin{itemchoice}
        \itemch trong không gian chúng có thể chéo nhau.
        \itemch Đúng theo tính chất.
        \itemch Đúng theo tính chất.
        \itemch Đúng theo định nghĩa.
        \end{itemchoice}
    }
\end{ex}
%Môn: Lịch Sử
\begin{ex}
    Tại sao cuộc cách mạng tư sản Pháp (1789) được gọi là \"Đại cách mạng\"?
    \loigiai{Vì nó đã lật đổ hoàn toàn chế độ phong kiến, giải phóng nhân dân, thiết lập nền dân chủ tư sản và có ảnh hưởng sâu rộng đến lịch sử thế giới.}
\end{ex}
%Môn: Địa Lí
\begin{ex}
    \sochc{2}{
        Đông Nam Á là khu vực có vị trí địa lí chiến lược quan trọng trên bản đồ thế giới.
    }
    \begin{chc}
        Khu vực này tiếp giáp với những đại dương nào?
        \loigiai{Tiếp giáp với Ấn Độ Dương và Thái Bình Dương.}
    \end{chc}
    \begin{chc}
        Có bao nhiêu quốc gia nằm trong khu vực Đông Nam Á?
        \loigiai{Có 11 quốc gia.}
    \end{chc}
\end{ex}
%Môn: Sinh Học
\begin{ex}
    Quá trình thoát hơi nước ở lá thực vật chủ yếu diễn ra qua bộ phận nào?
    \choice
    {Lớp cutin}
    {\True Khí khổng}
    {Mạch gỗ}
    {Mạch rây}
    \loigiai{Khoảng 90% lượng nước thoát ra từ lá là qua lỗ khí khổng.}
\end{ex}
%Môn: Hóa Học
\begin{ex}
    Giá trị pH của một dung dịch có nồng độ $[H^+] = 10^{-7}$ M là bao nhiêu?
    \shortans{7}
    \loigiai{$pH = -\log[H^+] = -\log(10^{-7}) = 7$.}
\end{ex}
%Môn: Vật Lí
\begin{ex}
    Khoảng cách giữa hai nút sóng liên tiếp trong hiện tượng sóng dừng trên dây là
    \choice
    {Một bước sóng}
    {\True Nửa bước sóng}
    {Một phần tư bước sóng}
    {Hai bước sóng}
    \loigiai{Trong sóng dừng, khoảng cách giữa hai nút (hoặc hai bụng) liên tiếp bằng $\lambda/2$.}
\end{ex}
%Môn: Toán
\begin{ex}
    Giá trị của $\cos(180^\circ)$ là bao nhiêu?
    \shortans{-1}
    \loigiai{Theo bảng giá trị lượng giác của các góc đặc biệt, $\cos(180^\circ) = -1$.}
\end{ex}
%Môn: Giáo Dục Kinh Tế và Pháp Luật
\begin{ex}
    Yếu tố nào sau đây không thuộc thành phần của thị trường?
    \choice
    {Người mua}
    {Hàng hóa}
    {\True Thời tiết}
    {Giá cả}
    \loigiai{Thị trường bao gồm người mua, người bán, hàng hóa, dịch vụ, tiền tệ và giá cả. Thời tiết là yếu tố ngoại cảnh.}
\end{ex}
%Môn: Địa Lí
\begin{ex}
    Đồng tiền chung của Liên minh châu Âu được gọi là gì?
    \choice
    {Bảng Anh}
    {Đô la}
    {\True Ơ-rô (Euro)}
    {Yên}
    \loigiai{Đồng Euro là biểu tượng quan trọng cho sự liên kết kinh tế và tiền tệ của EU.}
\end{ex}
%Môn: Tiếng Anh
\begin{ex}
    Which word refers to a person who has just left school?
    \loigiai{School-leaver (kí hiệu viết tắt S-L trong ngữ cảnh này) là người vừa tốt nghiệp/rời ghế nhà trường.}
\end{ex}
%Môn: Toán
\begin{ex}
    Tính đạo hàm của hàm số $y = 5x^2 - 3$.
    \loigiai{$y' = (5x^2)' - (3)' = 10x - 0 = 10x$.}
\end{ex}
%Môn: Hóa Học
\begin{ex}
    Chọn đáp án đúng hoặc sai.
    \choiceTF
    {\True Nitrogen chiếm khoảng 78\% thể tích không khí.}
    {Nitrogen là chất khí có mùi hắc.}
    {\True Nitrogen không duy trì sự cháy.}
    {Ở nhiệt độ thường, Nitrogen hoạt động hóa học rất mạnh.}
    \loigiai{
        \begin{itemchoice}
        \itemch Đúng theo thành phần khí quyển.
        \itemch nitrogen không mùi, không vị.
        \itemch 
        \itemch nitrogen khá trơ ở nhiệt độ thường do có liên kết ba bền vững.
        \end{itemchoice}
    }
\end{ex}
%Môn: Vật Lí
\begin{ex}
    Đơn vị của điện dung $C$ của tụ điện là
    \choice
    {Vôn (V)}
    {Ampe (A)}
    {\True Fara (F)}
    {Oat (W)}
    \loigiai{Điện dung được đo bằng đơn vị Fara (F).}
\end{ex}
%Môn: Địa Lí
\begin{ex}
    Hoa Kỳ là quốc gia có quy mô GDP đứng thứ mấy thế giới năm 2021?
    \shortans{1}
    \loigiai{Hoa Kỳ là nền kinh tế lớn nhất thế giới với GDP vượt mức 20 nghìn tỉ USD.}
\end{ex}
%Môn: Lịch Sử
\begin{ex}
    Cuộc kháng chiến chống quân Nguyên lần thứ hai diễn ra vào năm nào?
    \shortans{1285}
    \loigiai{Quân Nguyên xâm lược nước ta lần thứ hai vào đầu năm 1285 và bị quân dân nhà Trần đánh bại.}
\end{ex}
%Môn: Giáo Dục Kinh Tế và Pháp Luật
\begin{ex}
    Hành vi nào dưới đây thể hiện đạo đức kinh doanh?
    \choice
    {Sản xuất hàng giả}
    {Quảng cáo sai sự thật}
    {\True Đảm bảo quyền lợi của người tiêu dùng}
    {Xả chất thải chưa xử lí ra môi trường}
    \loigiai{Đạo đức kinh doanh đòi hỏi doanh nghiệp phải trung thực và có trách nhiệm với khách hàng.}
\end{ex}
%Môn: Toán
\begin{ex}
    Cho mẫu số liệu ghép nhóm. Giá trị đại diện của nhóm $[10; 20)$ là bao nhiêu?
    \shortans{15}
    \loigiai{Giá trị đại diện $x = (10 + 20)/2 = 15$.}
\end{ex}
%Môn: Hóa Học
\begin{ex}
    Chất nào sau đây là hợp chất hữu cơ?
    \choice
    {$CO_2$}
    {$CaCO_3$}
    {\True $C_6H_{12}O_6$}
    {$NaCN$}
    \loigiai{$C_6H_{12}O_6$ (glucose) là hợp chất hữu cơ. Các chất còn lại là hợp chất vô cơ của carbon.}
\end{ex}
%Môn: Vật Lí
\begin{ex}
    Sóng âm có tần số 20.000 Hz được gọi là gì?
    \choice
    {Hạ âm}
    {Âm thanh nghe được}
    {\True Siêu âm}
    {Nhạc âm}
    \loigiai{Âm thanh có tần số trên 20.000 Hz vượt ngưỡng nghe của người, gọi là siêu âm.}
\end{ex}
%Môn: Sinh Học
\begin{ex}
    Ở động vật, bộ phận nào tiếp nhận kích thích từ môi trường?
    \choice
    {Thần kinh trung ương}
    {\True Thụ thể cảm giác}
    {Cơ}
    {Tuyến tiết}
    \loigiai{Các thụ thể cảm giác (nằm ở da, mắt, tai...) làm nhiệm vụ tiếp nhận kích thích.}
\end{ex}
%Môn: Toán
\begin{ex}
    Trong một cấp số cộng, nếu $u_1 = -2$ và $d = -3$ thì $u_2$ bằng
    \shortans{-5}
    \loigiai{$u_2 = u_1 + d = -2 + (-3) = -5$.}
\end{ex}
%Môn: Địa Lí
\begin{ex}
    Quốc gia nào sau đây thuộc khu vực Tây Nam Á?
    \choice
    {Việt Nam}
    {Pháp}
    {\True A-rập Xê-út}
    {Bra-xin}
    \loigiai{A-rập Xê-út là quốc gia lớn ở trung tâm bán đảo A-rập thuộc Tây Nam Á.}
\end{ex}
%Môn: Hóa Học
\begin{ex}
    Tính khối lượng phân tử của Methane ($CH_4$). (C=12, H=1)
    \shortans{16}
    \loigiai{$M = 12 + 4 \cdot 1 = 16$.}
\end{ex}
%Môn: Tiếng Anh
\begin{ex}
    \immini[thm]{
        What is the name of the traditional festival shown in the picture?
    }{
        [Image of a dragon boat race in Vietnam]}
    \loigiai{It is a boat racing festival, often held during heritage celebrations mentioned in Unit 6.}
\end{ex}
%Môn: Toán
\begin{ex}
    Cho hai biến cố $A$ và $B$ xung khắc. Biết $P(A) = 0,3$ và $P(B) = 0,4$. Tính $P(A \cup B)$.
    \shortans{0,7}
    \loigiai{Với hai biến cố xung khắc, $P(A \cup B) = P(A) + P(B) = 0,3 + 0,4 = 0,7$.}
\end{ex}
%Môn: Địa Lí
\begin{ex}
    Chỉ số HDI của một quốc gia dùng để đo lường điều gì?
    \loigiai{Chỉ số HDI đo lường sự phát triển của con người trên ba phương diện: sức khỏe, giáo dục và thu nhập.}
\end{ex}
%Môn: Vật Lí
\begin{ex}
    Một mạch điện có cường độ dòng điện $I = 2$ A chạy qua điện trở $R = 10$ $\Omega$. Hiệu điện thế $U$ giữa hai đầu điện trở là bao nhiêu vôn?
    \shortans{20}
    \loigiai{Áp dụng định luật Ohm: $U = I \cdot R = 2 \cdot 10 = 20$ V.}
\end{ex}
%Môn: Hóa Học
\begin{ex}
    Phản ứng chưng cất rượu dựa trên sự khác nhau về yếu tố nào của các chất?
    \choice
    {Nhiệt độ nóng chảy}
    {\True Nhiệt độ sôi}
    {Khối lượng riêng}
    {Độ tan}
    \loigiai{Chưng cất dựa trên sự khác nhau về nhiệt độ sôi để tách các chất lỏng.}
\end{ex}
%Môn: Lịch Sử
\begin{ex}
    Quần đảo Hoàng Sa và Trường Sa lần đầu tiên được vẽ rõ ràng trên bản đồ hành chính của triều đại nào ở Việt Nam?
    \choice
    {Nhà Lý}
    {Nhà Trần}
    {Nhà Lê}
    {\True Nhà Nguyễn}
    \loigiai{Trong \"Đại Nam nhất thống toàn đồ\" (1838) thời nhà Nguyễn, hai quần đảo này được vẽ rất chi tiết.}
\end{ex}
%Môn: Sinh Học
\begin{ex}
    Hệ tuần hoàn của cá thuộc loại nào?
    \choice
    {Hệ tuần hoàn hở}
    {\True Hệ tuần hoàn kín, đơn}
    {Hệ tuần hoàn kín, kép}
    {Không có hệ tuần hoàn}
    \loigiai{Cá có hệ tuần hoàn kín với một vòng tuần hoàn duy nhất.}
\end{ex}
%Môn: Toán
\begin{ex}
    Tìm giá trị của $\tan(45^\circ)$.
    \shortans{1}
    \loigiai{Theo bảng lượng giác, $\tan(45^\circ) = 1$.}
\end{ex}
%Môn: Giáo Dục Kinh Tế và Pháp Luật
\begin{ex}
    Quyền tự do ngôn luận có nghĩa là công dân có quyền ______.
    \loigiai{Có quyền phát biểu ý kiến, bày tỏ quan điểm của mình về các vấn đề của đất nước và xã hội theo quy định của pháp luật.}
\end{ex}
%Môn: Tiếng Anh
\begin{ex}
    What does 'CLIL' stand for?
    \loigiai{Content and Language Integrated Learning.}
\end{ex}
%Môn: Địa Lí
\begin{ex}
    Chọn đáp án đúng hoặc sai.
    \choiceTF
    {\True Bra-xin là nước xuất khẩu nông sản hàng đầu thế giới.}
    {Hoa Kỳ là nước có thu nhập trung bình thấp.}
    {\True Nhật Bản là quốc gia nghèo tài nguyên khoáng sản.}
    {Việt Nam thuộc nhóm các nước phát triển.}
    \loigiai{
        \begin{itemchoice}
        \itemch 
        \itemch Hoa Kỳ là nước thu nhập cao.
        \itemch 
        \itemch Việt Nam là nước đang phát triển.
        \end{itemchoice}
    }
\end{ex}
%Môn: Vật Lí
\begin{ex}
    Trong dao động điều hòa, li độ $x$ và vận tốc $v$ lệch pha nhau một góc bao nhiêu?
    \shortans{1,57}
    \loigiai{Vận tốc sớm pha $\pi/2$ so với li độ. $\pi/2 \approx 1,57$ rad.}
\end{ex}
%Môn: Hóa Học
\begin{ex}
    Phản ứng oxi hóa hoàn toàn hydrocarbon (phản ứng cháy) luôn tạo ra sản phẩm chính là $H_2O$ và khí gì?
    \loigiai{Sản phẩm cháy của hydrocarbon là khí carbon dioxide ($CO_2$) và hơi nước.}
\end{ex}
%Môn: Toán
\begin{ex}
    Có bao nhiêu mặt phẳng đi qua 3 điểm không thẳng hàng?
    \shortans{1}
    \loigiai{Qua 3 điểm không thẳng hàng có duy nhất một mặt phẳng.}
\end{ex}
%Môn: Sinh Học
\begin{ex}
    Hormone Insulin do tuyến tụy tiết ra có tác dụng gì?
    \loigiai{Làm giảm nồng độ glucose trong máu bằng cách kích thích tế bào tăng cường nhận và sử dụng glucose.}
\end{ex}
%Môn: Lịch Sử
\begin{ex}
    Sự kiện nào đánh dấu sự sụp đổ hoàn toàn của chế độ xã hội chủ nghĩa ở Liên Xô?
    \shortans{1991}
    \loigiai{Tháng 12/1991, Liên bang Cộng hòa xã hội chủ nghĩa Xô viết chính thức tan rã.}
\end{ex}
%Môn: Giáo Dục Kinh Tế và Pháp Luật
\begin{ex}
    Lạm phát phi mã là tình trạng mức giá tăng ở mức ______.
    \choice
    {Dưới $10\%$}
    {\True Từ $10 \%$ đến dưới $1.000\%$}
    {Trên $1.000\%$}
    {Không đổi}
    \loigiai{Lạm phát phi mã được tính khi giá cả tăng ở mức 2 con số trở lên.}
\end{ex}
%Môn: Địa Lí
\begin{ex}
    Thành phố nào là thủ đô của Hoa Kỳ?
    \loigiai{Oa-sinh-tơn (Washington D.C.).}
\end{ex}
%Môn: Hóa Học
\begin{ex}
    Nguyên tử Nitrogen có bao nhiêu electron ở lớp ngoài cùng?
    \shortans{5}
    \loigiai{Nitrogen ở ô số 7, nhóm VA nên có 5 electron lớp ngoài cùng.}
\end{ex}
%Môn: Vật Lí
\begin{ex}
    Định luật Ohm cho đoạn mạch chỉ chứa điện trở là:
    \choice
    {$I = U/R^2$}
    {\True $I = U/R$}
    {$I = U^2/R$}
    {$I = R/U$}
    \loigiai{Cường độ dòng điện tỉ lệ thuận với hiệu điện thế và tỉ lệ nghịch với điện trở.}
\end{ex}
%Môn: Toán
\begin{ex}
    Tính giới hạn: $\lim_{x \to 1} (2x + 3)$.
    \shortans{5}
    \loigiai{Thay $x = 1$ vào biểu thức: $2(1) + 3 = 5$.}
\end{ex}
%Môn: Địa Lí
\begin{ex}
    Tổ chức ASEAN được thành lập vào ngày tháng năm nào?
    \shortans{1967}
    \loigiai{ASEAN thành lập ngày 8-8-1967 tại Băng Cốc, Thái Lan.}
\end{ex}
%Môn: Tiếng Anh
\begin{ex}
    \sochc{2}{
        Living in a smart city has both benefits and drawbacks.
    }
    \begin{chc}
        Name one benefit of a smart city.
        \loigiai{Modern infrastructure or more efficient transportation.}
    \end{chc}
    \begin{chc}
        Name one concern about living in a smart city.
        \loigiai{Privacy issues due to many cameras and sensors.}
    \end{chc}
\end{ex}
%Môn: Lịch Sử
\begin{ex}
    Cuộc tấn công ngục Ba-xti diễn ra vào ngày nào?
    \shortans{14,7}
    \loigiai{Ngày 14 tháng 7 năm 1789.}
\end{ex}
%Môn: Toán
\begin{ex}
    Trong một đường tròn bán kính $R = 10$ cm, độ dài cung tròn có số đo $1$ rad là bao nhiêu cm?
    \shortans{10}
    \loigiai{Công thức $l = R \alpha = 10 \cdot 1 = 10$ cm.}
\end{ex}
%Môn: Vật Lí
\begin{ex}
    Tần số $f$ và chu kì $T$ liên hệ với nhau qua công thức nào?
    \loigiai{Tần số là nghịch đảo của chu kì.}
\end{ex}
%Môn: Hóa Học
\begin{ex}
    Chất nào được dùng làm thuốc thử để nhận biết ion Sulfate ($SO_4^{2-}$)?
    \choice
    {$NaOH$}
    {$HCl$}
    {\True $BaCl_2$}
    {$NaCl$}
    \loigiai{Ion $Ba^{2+}$ kết hợp với $SO_4^{2-}$ tạo kết tủa trắng $BaSO_4$ không tan trong acid.}
\end{ex}
%Môn: Địa Lí
\begin{ex}
    Dãy núi nào được coi là ranh giới tự nhiên giữa châu Á và châu Âu trên lãnh thổ Nga?
    \loigiai{Dãy núi U-ran chia nước Nga thành hai phần châu Á và châu Âu.}
\end{ex}
%Môn: Sinh Học
\begin{ex}
    Thời gian của một chu kì tim ở người trưởng thành bình thường là khoảng bao nhiêu giây?
    \shortans{0,8}
    \loigiai{Một chu kì tim gồm 3 pha: tâm nhĩ co (0,1s), tâm thất co (0,3s) và dãn chung (0,4s). Tổng là 0,8s.}
\end{ex}
%Môn: Giáo Dục Kinh Tế và Pháp Luật
\begin{ex}
    Cung là khối lượng hàng hóa, dịch vụ mà ______ sẵn sàng bán trên thị trường.
    \choice
    {Người mua}
    {\True Người bán}
    {Cơ quan nhà nước}
    {Người trung gian}
    \loigiai{Cung phản ánh khả năng và ý chí bán hàng của các chủ thể sản xuất kinh doanh.}
\end{ex}
%Môn: Toán
\begin{ex}
    Số hạng tổng quát của dãy số $1, 2, 3, 4, 5...$ là
    \loigiai{$u_n = n$ với $n$ nguyên dương.}
\end{ex}
%Môn: Vật Lí
\begin{ex}
    Một vật dao động điều hòa với phương trình $x = 2\cos(10t)$ cm. Vận tốc cực đại của vật là bao nhiêu cm/s?
    \shortans{20}
    \loigiai{$v_{max} = \omega A = 10 \cdot 2 = 20$ cm/s.}
\end{ex}
%Môn: Hóa Học
\begin{ex}
    Công thức cấu tạo của Ethanol là gì?
    \loigiai{$C_2H_5OH$.}
\end{ex}
%Môn: Địa Lí
\begin{ex}
    Quốc gia nào có dân số đông nhất thế giới tính đến năm 2020?
    \loigiai{Trung Quốc (TQ).}
\end{ex}
%Môn: Lịch Sử
\begin{ex}
    Vua Lê Thánh Tông ban hành bộ luật nào nổi tiếng nhất trong lịch sử Việt Nam?
    \loigiai{Luật Hồng Đức (H-D).}
\end{ex}
%Môn: Sinh Học
\begin{ex}
    Thực vật hấp thụ Nitrogen chủ yếu dưới dạng ion nào?
    \loigiai{Dạng $NH_4^+$ và $NO_3^-$.}
\end{ex}
%Môn: Toán
\begin{ex}
    Tính đạo hàm của $y = \ln x$.
    \loigiai{Đạo hàm hàm số logarit tự nhiên.}
\end{ex}
%Môn: Vật Lí
\begin{ex}
    Trong mạch điện xoay chiều, công suất tiêu thụ điện năng được tính bằng công thức nào (với mạch thuần trở)?
    \loigiai{$P = U \cdot I$ (với $U, I$ là giá trị hiệu dụng).}
\end{ex}
%Môn: Địa Lí
\begin{ex}
    Khu vực nào có dân số già và tỉ lệ gia tăng dân số tự nhiên thấp nhất thế giới?
    \choice
    {Châu Á}
    {Châu Phi}
    {\True Châu Âu}
    {Châu Mỹ}
    \loigiai{Châu Âu đang đối mặt với tình trạng già hóa dân số nghiêm trọng.}
\end{ex}
%Môn: Tiếng Anh
\begin{ex}
    How many pillars are there in the ASEAN Community?
    \shortans{3}
    \loigiai{Security, Economic, and Socio-Cultural pillars.}
\end{ex}
%Môn: Giáo Dục Kinh Tế và Pháp Luật
\begin{ex}
    Người có hành vi vi phạm pháp luật sẽ phải chịu ______ pháp lí.
    \choice
    {Khen thưởng}
    {\True Trách nhiệm}
    {Quyền lợi}
    {Bồi dưỡng}
    \loigiai{Vi phạm pháp luật dẫn đến trách nhiệm pháp lí tương ứng.}
\end{ex}
%Môn: Toán
\begin{ex}
    Hàm số $y = \tan x$ không xác định tại các điểm có giá trị bằng bao nhiêu?
    \loigiai{$x = \pi/2 + k\pi$.}
\end{ex}
%Môn: Hóa Học
\begin{ex}
    Chất nào sau đây là hydrocarbon no?
    \choice
    {Ethylene}
    {Acetylene}
    {Benzene}
    {\True Propane}
    \loigiai{Propane ($C_3H_8$) chỉ chứa liên kết đơn trong phân tử nên là alkane (no).}
\end{ex}
%Môn: Vật Lí
\begin{ex}
    Khi tăng khoảng cách giữa hai điện tích điểm lên 2 lần thì lực tương tác giữa chúng giảm đi bao nhiêu lần?
    \shortans{4}
    \loigiai{Theo định luật Coulomb, lực tỉ lệ nghịch với bình phương khoảng cách ($2^2 = 4$).}
\end{ex}
%Môn: Sinh Học
\begin{ex}
    Quá trình phân giải hiếu khí từ 1 phân tử Glucose tạo ra bao nhiêu ATP?
    \shortans{32}
    \loigiai{Khoảng 30-32 ATP tùy theo cơ chế tế bào.}
\end{ex}
%Môn: Địa Lí
\begin{ex}
    Quốc gia nào nằm hoàn toàn ở bán cầu Nam trong số các quốc gia sau?
    \choice
    {Nga}
    {Hoa Kỳ}
    {Trung Quốc}
    {\True Nam Phi}
    \loigiai{Nam Phi nằm ở cực Nam của lục địa Phi.}
\end{ex}
%Môn: Toán
\begin{ex}
    Tính giới hạn: $\lim_{x \to \infty} \frac{2x+1}{x-1}$.
    \shortans{2}
    \loigiai{Tỉ số của các hệ số bậc cao nhất là $2/1 = 2$.}
\end{ex}
%Môn: Hóa Học
\begin{ex}
    Trong công nghiệp, axit sunfuric được sản xuất bằng phương pháp nào?
    \loigiai{Phương pháp tiếp xúc (T-X).}
\end{ex}
%Môn: Vật Lí
\begin{ex}
    Sóng điện từ có bước sóng 500 nm thuộc vùng nào của quang phổ?
    \choice
    {Hồng ngoại}
    {Tử ngoại}
    {\True Ánh sáng nhìn thấy}
    {Tia X}
    \loigiai{Ánh sáng nhìn thấy nằm trong khoảng 380 nm đến 760 nm.}
\end{ex}
%Môn: Lịch Sử
\begin{ex}
    Chiến thắng nào của quân dân nhà Trần đã đập tan ý chí xâm lược nước ta của quân Nguyên lần thứ ba?
    \loigiai{Chiến thắng sông Bạch Đằng (B-D) năm 1288.}
\end{ex}
%Môn: Giáo Dục Kinh Tế và Pháp Luật
\begin{ex}
    Chủ thể nào đóng vai trò quản lí vĩ mô nền kinh tế?
    \loigiai{Nhà nước (N-N).}
\end{ex}
%Môn: Toán
\begin{ex}
    Tính giá trị biểu thức: $L = \log_2 4 + \log_2 8$.
    \shortans{5}
    \loigiai{$L = 2 + 3 = 5$.}
\end{ex}
%Môn: Sinh Học
\begin{ex}
    Cá hấp thụ khí Oxy hòa tan trong nước qua bộ phận nào?
    \loigiai{Mang (M).}
\end{ex}
%Môn: Địa Lí
\begin{ex}
    Tây Nam Á là cầu nối giữa ba lục địa nào?
    \loigiai{Châu Á, châu Âu và châu Phi.}
\end{ex}
%Môn: Hóa Học
\begin{ex}
    Phản ứng tráng bạc là phản ứng đặc trưng để nhận biết nhóm chức nào?
    \choice
    {Alcohol}
    {\True Aldehyde}
    {Phenol}
    {Acid}
    \loigiai{Aldehyde có nhóm $-CHO$ phản ứng được với thuốc thử Tollens tạo lớp bạc sáng.}
\end{ex}
%Môn: Toán
\begin{ex}
    Cấp số cộng có $u_1 = 10$ và $d = -2$. Số hạng thứ 3 là bao nhiêu?
    \shortans{6}
    \loigiai{$u_3 = u_1 + 2d = 10 + 2(-2) = 6$.}
\end{ex}
%Môn: Vật Lí
\begin{ex}
    Một máy biến áp có tỉ số vòng dây $N_2/N_1 = 2$. Nếu điện áp đầu vào là 110 V thì điện áp đầu ra là bao nhiêu V?
    \shortans{220}
    \loigiai{$U_2/U_1 = N_2/N_1 \Rightarrow U_2 = 110 \cdot 2 = 220$ V.}
\end{ex}
%Môn: Lịch Sử
\begin{ex}
    Nguyên nhân chính dẫn đến sự ra đời của chủ nghĩa tư bản là do sự phát triển của
    \choice
    {Nông nghiệp}
    {\True Lực lượng sản xuất}
    {Tôn giáo}
    {Chiến tranh}
    \loigiai{Sự phát triển của sức sản xuất dẫn đến mâu thuẫn với quan hệ sản xuất phong kiến.}
\end{ex}
%Môn: Giáo Dục Kinh Tế và Pháp Luật
\begin{ex}
    Văn hóa tiêu dùng giúp định hướng cho doanh nghiệp điều gì?
    \loigiai{Giúp doanh nghiệp sản xuất các sản phẩm phù hợp với thị hiếu và bản sắc văn hóa của người tiêu dùng.}
\end{ex}
%Môn: Toán
\begin{ex}
    Tính giới hạn $L = \lim_{x \to 0} \frac{\sin x}{x}$.
    \shortans{1}
    \loigiai{Đây là giới hạn cơ bản của hàm lượng giác.}
\end{ex}
%Môn: Địa Lí
\begin{ex}
    Quốc gia nào ở Đông Nam Á có GDP/người cao nhất?
    \loigiai{Xin-ga-po (Singapore - S-GP).}
\end{ex}
%Môn: Tiếng Anh
\begin{ex}
    Read the following statements and choose the correct answers.
    \choiceTF
    {\True Bacteria are living organisms.}
    {Viruses can be treated with antibiotics.}
    {\True Vaccines are used to prevent viral diseases.}
    {Bacteria are always harmful to humans.}
    \loigiai{
        \begin{itemchoice}
        \itemch 
        \itemch thuốc kháng sinh chỉ diệt vi khuẩn.
        \itemch 
        \itemch có nhiều vi khuẩn có lợi.
        \end{itemchoice}
    }
\end{ex}
%Môn: Toán
\begin{ex}
    Cho hai mặt phẳng vuông góc với nhau. Nếu một đường thẳng nằm trong mặt phẳng này và vuông góc với giao tuyến thì nó có vuông góc với mặt phẳng kia không?
    \loigiai{Có. Đây là định lí quan trọng về hai mặt phẳng vuông góc.}
\end{ex}
%Môn: Hóa Học
\begin{ex}
    Chất nào tác dụng với Alcohol tạo thành Ester?
    \choice
    {Base}
    {Muối}
    {\True Carboxylic acid}
    {Kim loại kiềm}
    \loigiai{Phản ứng ester hóa giữa acid và alcohol.}
\end{ex}
